\documentclass[10pt,a4paper]{amsart}
\usepackage[margin=19mm]{geometry}
\usepackage{amsmath,amssymb,mathtools}
\usepackage[scr=rsfs,cal=euler]{mathalfa}
\usepackage{xfrac}
\usepackage[unicode,hidelinks,bookmarks=false]{hyperref}
\newcommand{\ie}{i.e.\ }
\newcommand{\cf}{cf.\ }
\newcommand{\resp}{respectively\ }
\newcommand{\Subsection}{\subsection}
\providecommand{\nobreakdash}{\nobreak\hbox{-}}
\setlength{\emergencystretch}{2em}
\multlinegap0pt
\usepackage{tikz}
\usepackage{amssymb}
\usepackage{dsfont}
\usepackage{enumerate}
\usepackage[shortlabels]{enumitem}
\newtheorem{lemma}{Lemma}
\newtheorem{theorem}{Theorem}
\title{Furtherness in finite topological spaces}
\author{Akhilesh Badra, Hemant Kumar Singh}
\date{Source: arXiv:2603.18580v1 (CC BY 4.0)}
\begin{document}
\maketitle

\noindent\emph{Source and licence.} Furtherness in finite topological spaces. Version arXiv:2603.18580v1 (CC BY 4.0), distributed by arXiv under the Creative Commons Attribution 4.0 International licence, http://creativecommons.org/licenses/by/4.0/, as recorded in the arXiv licence metadata for this exact version and shown on the abstract page https://arxiv.org/abs/2603.18580v1. Full licence notice: \texttt{licenses/arxiv-2603.18580-CC-BY-4.0.txt}.\par\smallskip

\noindent\textit{Editorial scope.} Counted unit: the article's theorem 2.9. (source0.tex lines 321--323). The article's complete original proof of it is delivered with it (source0.tex lines 324--346, one \texttt{proof} environment of 23 lines). No mathematics is written, repaired or re-derived: the statement and the proof are the article's own lines, in the article's own notation, and no retained line is substituted or re-flowed.  Retained uncounted as the source-local context the proof needs: lines 309--311.  Nothing was omitted from any retained span, so the package carries no omission record.  No retained line contains a citation, so the wrapper carries no bibliography and no bibliography entry.  No retained line contains a figure environment, an image-inclusion command or drawing code, so no image is delivered and the package declares no asset dependency.  Editorial formatting only, in addition: an AMS \texttt{amsart} preamble that restates the article's own declarations and macros used by the retained lines, namely the environments \texttt{lemma}, \texttt{theorem} on separate counters, together with the standard packages those lines need.  The retained environments print in this document as Lemma 1, Theorem 1 in order of appearance, whereas the article prints the same items with the numbering of its own full document.  The complete native source of the article as distributed in the arXiv e-print for this exact version is bundled unchanged as \texttt{sources/196706/source0.tex}.
\begin{lemma}\label{2.8.}
	Let $X$ be a finite space and $(U_j)_{j\geq0}$ be a nested sequence of open sets around $x\in X$, then $U_j= U_{j-1}\cup U_a$, for some $a\in X\backslash U_{j-1}$.
\end{lemma}

\begin{theorem}\label{2.9.}
	Let $X$ be a finite space and $x,y\in X$ such that $\Psi(x,y)=k$. Then there exists a nested sequence $(U_j)_{j\geq 0}$ of open sets around $x$ such that $U_k=U_{x,y}$.
\end{theorem}

\begin{proof} 
	As  $\Psi(x,y)=k$, there exists a nested sequence $(V_j)_{j\geq 0}$ of open sets around $x$ such that $y\in V_k$. If $k=0$, then it is trivialy true because $y\in V_0=U_x\implies U_y\subseteq U_x=U_{x,y}=V_0$. Now let $k>0$. By Lemma \ref{2.8.}, we have $V_0= U_x$, $V_1= V_0\cup U_{a_1},$ for some $a_1\notin V_0$, $V_2= V_0\cup U_{a_1}\cup U_{a_2},$ for some  $a_2 \notin V_1$. By induction, we get $V_{k-1}= V_0\cup (\bigcup\limits_{i=1}^{k-1} U_{a_i})$, and $V_k= V_{k-1}\cup U_{a_k}$ for some $a_k\notin V_{k-1}$. As $y\notin V_{k-1}\implies y\in U_{a_k}$. So, $U_y\subseteq U_{a_k} \implies V_{k-1}\subset V_{k-1}\cup U_y\subseteq V_{k-1}\cup  U_{a_k} =V_k.$ %\implies  V_{k-1}\cup U_y= V_{k-1}\cup  U_{a_k} $.
	So, $V_k= V_{k-1}\cup U_y$. \par 
	As, $U_x$ is contained in open set $U_{x,y},$ there exists a nested sequence $(W_j)_{j\geq 0}$ around $x$ such that $W_p=U_{x,y}$ for some $p\in \mathbb{N}\cup\{0\}$. Now, $y\in W_p\implies p\geq k.$ Again by lemma \ref{2.8.}, we have $W_0= U_x$, and
	% $W_1= U_x\cup U_{b_1},$ for some $b_1\notin W_0$, $W_2= U_x\cup U_{b_1}\cup U_{b_2},$ for some  $b_2 \notin W_1$, similarly,
	$W_p= U_x\cup (\bigcup\limits_{j=1}^p U_{b_j})$, for $b_j \notin W_{j-1}$.
	% So, $U_x\cup U_y =  U_x\cup (\bigcup\limits_{j=1}^p U_{b_j})$.
	
	
	% So, $U_{b_j}\subseteq U_y,  1\leq j\leq p$. Therefore, $\bigcup\limits_{i=1}^p U_{b_i}\subseteq U_y$. As $k>0$ and $y\notin U_x$, and $y\in U_x\cup U_y =  U_x\cup (\bigcup\limits_{i=1}^p U_{b_i})$. So, $y\in \bigcup\limits_{i=1}^p U_{b_i}$. Therefore, $U_y\subseteq\bigcup\limits_{j=1}^p U_{b_j}$. So, $U_y=\bigcup\limits_{j=1}^p U_{b_j}$.\par 
	
	% As $x,y\in V_k$, and $W_p=U_{x,y}\implies W_p\subseteq V_k.$ So, $ U_x \cup (\bigcup\limits_{i=1}^p U_{b_i})\subseteq  U_x\cup (\bigcup\limits_{j=1}^{k} U_{a_j}) \implies \bigcup\limits_{i=1}^p U_{b_i}\subseteq \bigcup\limits_{j=1}^{k} U_{a_j}.$ Which means each $b_j$ belongs to some $U_{a_i}$, for $ 1\leq i\leq k, 1\leq j\leq p$. That implies each $b_j$ belongs to some $V_i$, for $ 1\leq i\leq k, 1\leq j\leq p$.\par
	%Now, we observe that $b_j \in U_{a_i} $, for some $ 1\leq i\leq k-1,~ 1\leq j\leq p$. If there exist a $b_j$ such that $b_j\notin U_{a_i},\forall 1\leq i\leq (k-1) \implies U_{b_j}\nsubseteq V_{k-1}$. Then $V_{k-1}\subseteq V_{k-1}\cup U_{b_j}\subseteq  V_{k}$. That means $ V_{k-1}\cup U_{b_j}$ is an open subset betweeb $ V_{k-1}$ and $ V_{k}$, a contradiction. So, $b_j$ must belongs to some $U_{a_i},  \forall ~ 1\leq j\leq p$.\par 
	
	% We already have $V_{k-1}= U_x\cup (\bigcup\limits_{i=1}^{k-1} U_{a_i})$. So, we get  $V_k= V_{k-1}\cup U_y=  U_x\cup (\bigcup\limits_{i=1}^{k-1} U_{a_i})\cup (\bigcup\limits_{j=1}^p U_{b_j})$ $\implies$
	Now, $b_j\in  U_x \cup (\bigcup\limits_{j=1}^p U_{b_j})= U_x\cup U_y$ and $b_j\notin U_x \implies b_j\in U_y,  1\leq j\leq p$.
	So, $b_j\in V_k, 1\leq j\leq p.$ Next, we observe that each  $V_l$ contains atmost one $b_j$ which is not in $V_{l-1}, 1\leq l\leq k$. 
	%Let $b_j, b_{j'}\in U_{a_l} $ for some $1\leq l\leq (k-1)$ and $1\leq j,j' \leq p$. Which means $U_{b_j}\subseteq U_{a_l}$ and $U_{b_{j'}}\subseteq U_{a_l}$. As $V_{l-1}=  U_x \bigcup\limits_{i=1}^{l-1} U_{a_i}$
	%Let $b_j\cup b_{j'}\subseteq V_1, $ for $1\leq j,j' \leq p$. Then $ V_{0}=U_x\subset U_x\cup U_{b_j}\subset U_x\cup U_{b_j}\cup U_{b_{j'}} \subseteq V_{1}$, a contradiction. Similarly, i
	If $b_j, b_{j'}\in V_l $ for some $1\leq l\leq k$ such that $b_j, b_{j'} \notin V_{l-1}$ and $b_j\neq b_{j'}$. Then $ V_{l-1}\subset  V_{l-1}\cup U_{b_j}\subset V_{l-1}\cup U_{b_j} \cup  U_{b_{j'}} \subseteq V_{l}$, which contradicts that there is no open set between $V_{l-1}$ and $V_l$. So, $V_k$ can have atmost $k$ distinct $b_j$ for $1\leq j\leq p$. But $b_j\in V_k,\forall~ 1\leq j\leq p.$ So, $p\leq k$.
	%If $p>k$, then by peigeon hole principle there exist at least one $U_{a_i}$ for $1\leq i \leq (k-1)$ such that $U_{a_i}$ contains more than one $b_j$ which are not in $U_{a_{i-1}}$ for $1\leq j\leq p$. Which is not true.
	Thus, $p=k$. So, there exist a nested sequence $(U_j)_{j\geq 0}$ of open sets around $x$ such that $U_k=U_{x,y}$.
\end{proof}

\end{document}
